% 修士論文 本体(lualatex でビルドする: latexmk -lualatex main.tex、または lualatex main.tex を 2 回)
\documentclass[11pt,a4paper,openany]{ltjsreport}
\usepackage[haranoaji]{luatexja-preset}
\usepackage{amsmath,amssymb,amsthm,bm}
\usepackage{graphicx}
\usepackage{booktabs}
\usepackage{array}
\usepackage[margin=25mm]{geometry}
\usepackage{caption}
\usepackage{xcolor}
\usepackage[hidelinks]{hyperref}

\graphicspath{{../}{figures/}}
\captionsetup{font=small,labelfont=bf}
\allowdisplaybreaks

% ---- 記号 ----
\newcommand{\bra}[1]{\langle #1|}
\newcommand{\ket}[1]{|#1\rangle}
\newcommand{\ev}[1]{\langle #1\rangle}
\newcommand{\Tr}{\operatorname{Tr}}
\newcommand{\Var}{\operatorname{Var}}
\newcommand{\Cov}{\operatorname{Cov}}
\newcommand{\cd}{c^\dagger}
\newcommand{\up}{\uparrow}
\newcommand{\dn}{\downarrow}
\newcommand{\Sv}{\mathbf S}

% ---- 執筆中の印(提出前にすべて消す) ----
\newcommand{\todo}[1]{\textcolor{red}{\textbf{[TODO: #1]}}}

\newtheorem{proposition}{命題}[chapter]

\title{\bfseries 共通ランダム測定によるハバード模型の局所観測量の推定\\[4pt]
\large — 古典計算による近似状態(prior)の選び方と、その物理 —}
\author{龍山 昊平}
\date{\todo{提出年月}}

\begin{document}
\maketitle

\begin{abstract}
古典シャドウ(ランダムな基底での測定)に、古典計算で作った近似状態(prior)を組み合わせて推定の分散を減らす共通ランダム測定(CRM)\cite{vermersch2024}を、フェルミオンのハバード模型に系統的に適用した。単一のパウリ文字列に対する CRM の分散は元論文が厳密に与えており、分散を減らせる割合(利得)は prior の大域的な忠実度ではなく、その観測量に対する prior の外れ $\Delta$ で決まる。本研究は、ハバード模型の現実的な prior について、この外れが観測量ごとにどう振る舞い、なぜそうなるかを調べた。

第一に、系を 256 量子ビットまで広げて prior の大域的な忠実度が崩れても、局所観測量の利得はほとんど変わらないことを示した。第二に、平均場(非制限 Hartree--Fock)prior は、電荷の量(オンサイトの ZZ 相関)では強結合で誤差が $(t/U)^2$ で消えて利得の天井に届く一方、スピンの量では損をすることを示した。その原因を、スピンの向きを選んでしまうことと、隣どうしの量子的な一重項の相関が足りないことの2つに分けた。前者はスピン回転で平均した混合状態 prior で直り、後者は直らない。第三に、スピン射影 HF の効果が系の大きさに反比例して消えることを回転子の描像で定量的に説明した。結合次元を制限した変分 MPS をスピン回転で平均した prior は、電荷とスピンの両方で損をせず、厳密な状態を使わない prior の中で最も良いことを示した。第四に、二重占有のようなパウリ項の和について CRM の分散を厳密に書き下し、利得が項ごとの外れで決まることを示した。さらに、2 次元シリンダー(最大 96 量子ビット)とドープ系での平均場 prior の有効範囲、ショットの配分、損をしない推定量、読み出し誤差に対する脆さと prior 側での補正を調べた。\todo{2次元の再計算の結果が出たら数値を入れて要旨を確定する}

\end{abstract}

\tableofcontents

% =====================================================================
\chapter{序論}
\label{ch:intro}
% =====================================================================

\section{背景:量子多体系の性質を測るコスト}

量子計算機や量子シミュレータ(超伝導量子ビット、イオントラップ、冷却原子など)で多体の量子状態を用意できるようになり、その状態の性質を\emph{測定で取り出す}ことが中心的な課題になっている。量子力学の測定は確率的なので、ある物理量の期待値を精度 $\epsilon$ で知るには、おおよそ $1/\epsilon^2$ 回の測定を繰り返す必要がある。系が大きくなり、知りたい物理量が多くなるほど、この測定の回数(ショット数)と、測定のたびに回路を組み替える回数(測定の\emph{設定}の数)が実験の実質的なコストになる。

この問題に対する有力な手法が\emph{古典シャドウ}(classical shadow)\cite{huang2020}である。各量子ビットを独立にランダムな基底($X,Y,Z$ のいずれか)で測り、その結果から状態の「影」を古典計算機上に作る。1組の測定データから多数の局所的な物理量を同時に推定できることが利点である。一方で、観測量の台(作用する量子ビット)が $\lvert A\rvert$ 個のとき、推定の分散は $3^{\lvert A\rvert}$ に比例して増える。これは、ランダムに選んだ基底が観測量の台の上でたまたま全部一致したときにしか、その観測量の情報が得られないためである。

Vermersch らは、この古典シャドウに古典計算で作った近似状態を組み合わせる\emph{共通ランダム測定}(common randomized measurements, CRM)を提案した\cite{vermersch2024}。実験の影から「同じ基底で近似状態を測ったときの影」を差し引き、近似状態そのものを足し戻す。近似状態(以下 \emph{prior} と呼ぶ)が何であっても推定は偏らず、prior が良いほど分散が減る。統計学でいう\emph{共通乱数法}・\emph{制御変量法}の量子版である。

\section{本研究の問い}

CRM の効果は prior の良さで決まる。では、量子多体系、とくに強相関電子系の典型例であるハバード模型に対して、
\begin{enumerate}
  \item \textbf{どんな prior が、どの観測量で、どれだけ効くのか。}
  \item \textbf{それはなぜか} — prior のどの物理が、どの観測量の誤差に効いているのか。
\end{enumerate}
これが本研究の問いである。

元論文は、単一のパウリ文字列に対する CRM の分散を厳密に与えている(arXiv 版の付録 B.1 の式(25)、PRX Quantum 版の付録 C.1 の式(C7))。その式によれば、分散を決めるのは prior の\emph{大域的な忠実度}ではなく、その観測量に対する prior の外れ
\[
  \Delta=\Tr[P\rho]-\Tr[P\sigma]
\]
である(第\ref{ch:background}章)。したがって上の問いは、「ハバード模型の現実的な prior は、どの観測量でどれだけ外れるのか、それはなぜか」という\emph{物理の問い}に置き換わる。本研究の主な寄与はこの置き換えた問いに答えたことにある。

\section{本研究の寄与}

本研究で明らかにしたことを、元論文にあるものと区別して挙げる。

\begin{description}
  \item[元論文にあるもの(本研究は実装で確かめた)] 単一パウリ文字列の分散の厳密な式、その比として出る利得の式、損得の境目 $\lvert\Delta\rvert\le\lvert\Tr[P\rho]\rvert$、prior によらない不偏性。
  \item[本研究の寄与]
\end{description}
\begin{enumerate}
  \item \textbf{局所性}(第\ref{ch:locality}章):系を $L=8$ から $128$ サイト(256 量子ビット)に広げると prior の大域的な忠実度は崩れるが、局所観測量の利得はほとんど変わらないことを示した。prior の良さの指標としては、観測量の台に縮約したトレース距離が $\lvert\Delta\rvert$ の最良の上界で、半充填では等号が成り立つことも示した。
  \item \textbf{観測量ごとに prior の得意・不得意が逆転すること}(第\ref{ch:priors}章):平均場(非制限 Hartree--Fock, UHF)prior は電荷の量(二重占有に関わる量)では強結合で誤差が $(t/U)^2$ で消える一方、スピンの量では損をする。その原因を2つ(スピンの向きを選ぶこと、量子的な一重項の相関が足りないこと)に分け、それぞれの直し方と限界を示した。
  \item \textbf{prior の比較}(第\ref{ch:priors}章):スピン回転で平均した UHF(UHF-sym)、スピン射影 HF(PHF)、切断 MPS、変分 MPS、変分 MPS の回転平均を同じ厳密な参照状態で比べた。PHF の効果が $1/L$ で消えることを回転子の描像で定量的に説明した。変分 MPS を回転平均した prior が、厳密な状態を使わない prior の中で最も良いことを示した。
  \item \textbf{和の観測量の厳密な分散}(第\ref{ch:methods}章、付録\ref{app:sumvar}):二重占有のように複数のパウリ項の和で書かれる観測量について、CRM の分散を厳密に書き下した。利得は合計の $\Delta$ ではなく項ごとの外れで決まる。
  \item \textbf{2 次元とドープ系}(第\ref{ch:2d}章):最大 96 量子ビットの 2 次元シリンダーで、平均場 prior が MPS prior に代われる範囲と、ドープ系の電荷の模様で破綻する場所を調べた。
  \item \textbf{実用上の問題}(第\ref{ch:practical}章):ショットの配分、損をしない推定量(最適係数)、読み出し誤差に対する脆さと prior 側での補正、matchgate シャドウとの関係。
\end{enumerate}

\section{本論文の構成}

第\ref{ch:background}章でハバード模型・古典シャドウ・CRM を導入し、元論文の分散の式をまとめる。第\ref{ch:methods}章で本研究の手法(和の観測量の厳密な分散、prior の作り方、参照状態と数値実験)を述べる。第\ref{ch:locality}章から第\ref{ch:practical}章が結果である。第\ref{ch:conclusion}章で結論と今後の課題を述べる。付録に、半充填でサイトごとの電子数が厳密に 1 になる理由(粒子正孔対称性)、和の観測量の分散の導出、計算の再現方法をまとめた。

% =====================================================================
\chapter{背景}
\label{ch:background}
% =====================================================================

\section{ハバード模型}
\label{sec:hubbard}

\subsection{模型}
本研究で扱うのは、格子上の電子のハバード模型
\begin{equation}
  H=-t\sum_{\langle ij\rangle,\sigma}\bigl(\cd_{i\sigma}c_{j\sigma}+\mathrm{h.c.}\bigr)+U\sum_i n_{i\up}n_{i\dn}-\mu\sum_{i}n_i
  \label{eq:hubbard}
\end{equation}
である。$c_{i\sigma}$ はサイト $i$、スピン $\sigma\in\{\up,\dn\}$ の電子の消滅演算子、$n_{i\sigma}=\cd_{i\sigma}c_{i\sigma}$、$n_i=n_{i\up}+n_{i\dn}$ で、$\langle ij\rangle$ は隣接するサイトの組を表す。第 1 項は電子が隣のサイトへ飛び移る運動エネルギー、第 2 項は同じサイトに 2 個の電子がいるときのクーロン反発である。全体を通じて $t=1$ とする。$\mu=U/2$ のとき、二部格子上で粒子数がサイト数に等しい\emph{半充填}になる。

格子は 1 次元の鎖(開放端・周期境界)と、幅 $W$ のシリンダー(円周方向に周期、軸方向に開放)、および小さな 2 次元トーラスを使う。

\subsection{強結合の極限と超交換}
半充填で $U\gg t$ のとき、二重占有はエネルギー $U$ を要するので抑えられ、各サイトにほぼ 1 個の電子が局在する(Mott 絶縁体)。残ったスピンの自由度の間には、2 次の摂動で
\begin{equation}
  H_{\rm eff}=J\sum_{\langle ij\rangle}\Bigl(\Sv_i\cdot\Sv_j-\tfrac14n_in_j\Bigr),\qquad J=\frac{4t^2}{U}
  \label{eq:heisenberg}
\end{equation}
の反強磁性的な相互作用(超交換)が働く。Hellmann--Feynman の定理 $\partial E_0/\partial U=\sum_i\ev{n_{i\up}n_{i\dn}}$ を式\eqref{eq:heisenberg}のエネルギーに当てはめると、強結合での二重占有率は
\begin{equation}
  d\equiv\ev{n_{i\up}n_{i\dn}}\simeq\frac{4t^2}{U^2}\sum_{j\in\partial i}\frac12\Bigl(\frac14-\ev{\Sv_i\cdot\Sv_j}\Bigr)
  \label{eq:d-strong}
\end{equation}
となる($\partial i$ はサイト $i$ の隣、1 次元では $\tfrac12\sum_{j\in\partial i}$ が 1 本分の結合になる)。\textbf{二重占有は隣どうしのスピンがどれだけ一重項を組んでいるかを測っている}。この関係は第\ref{ch:priors}章で平均場 prior の誤差を解析するときに中心的な役割を果たす。

\subsection{Jordan--Wigner 変換}
\label{sec:jw}
電子の模型を量子ビットで扱うため、軌道(サイトとスピンの組)に一列に番号を振り、Jordan--Wigner 変換
\begin{equation}
  c_q=\Bigl(\prod_{k<q}Z_k\Bigr)\frac{X_q+iY_q}{2}
\end{equation}
で写す。占有数は $n_q=(1-Z_q)/2$、すなわち $Z_q=1-2n_q$ である。本研究では各サイトの $\up$ と $\dn$ を隣り合う量子ビットに置き、2 次元格子はシリンダーの円周方向に沿って蛇行するように並べる。

この写像のもとで、密度やスピンの $z$ 成分の観測量は $Z$ だけの短いパウリ文字列になる。例えば
\begin{equation}
  Z_{i\up}Z_{i\dn}=1-2n_i+4n_{i\up}n_{i\dn},\qquad
  n_{i\up}n_{i\dn}=\tfrac14\bigl(1-Z_{i\up}-Z_{i\dn}+Z_{i\up}Z_{i\dn}\bigr)
  \label{eq:docc-pauli}
\end{equation}
である。半充填では $\ev{n_i}=1$ なので(付録\ref{app:ph})、オンサイトの $ZZ$ 相関と二重占有率は $\ev{Z_{i\up}Z_{i\dn}}=4d-1$ で 1 対 1 に対応する。一方ホッピングは、2 つの軌道の間にある量子ビットの $Z$ の列(JW 弦)を伴い、台が長くなる。

\section{古典シャドウ}
\label{sec:shadow}

\subsection{プロトコル}
$N$ 量子ビットの状態 $\rho$ の各量子ビットについて、基底 $b_q\in\{X,Y,Z\}$ を独立に一様に選んで測る。基底の組 $u=(b_1,\dots,b_N)$ を\emph{設定}と呼ぶ。設定を $n_u$ 通り選び、それぞれで $n_m$ 回測る。

台 $A$($\lvert A\rvert$ 個の量子ビット)のパウリ文字列 $P=\bigotimes_{q\in A}\sigma_{a_q}$ の推定値は、1 つの設定 $u$ と測定結果 $s_q\in\{\pm1\}$ について
\begin{equation}
  X(P;u,s)=3^{\lvert A\rvert}\Bigl(\prod_{q\in A}s_q\Bigr)\,\mathbf 1[\forall q\in A:\ b_q=a_q]
  \label{eq:snapshot}
\end{equation}
である。基底が台の上ですべて一致したとき(確率 $3^{-\lvert A\rvert}$)だけ $3^{\lvert A\rvert}$ 倍された値を返し、それ以外は $0$ を返す。

\subsection{分散}
1 つの設定での $n_m$ ショットの平均を $\bar X(u)$ と書く。基底が一致した設定では $\prod_q s_q$ は平均 $x\equiv\Tr[P\rho]$、分散 $1-x^2$ の $\pm1$ の確率変数である。全分散の公式で、設定の選び方による揺らぎとショットの揺らぎに分けると
\begin{equation}
  n_u\Var[\text{標準}]=\underbrace{(3^{\lvert A\rvert}-1)x^2}_{\text{基底のくじ}}+\underbrace{\frac{3^{\lvert A\rvert}(1-x^2)}{n_m}}_{\text{ショット雑音 }v_s}
  \label{eq:var-std}
\end{equation}
となる。第 1 項は、基底が当たるか外れるかで推定値が $3^{\lvert A\rvert}x$ と $0$ の間を行き来することによる。これが古典シャドウの分散の大部分を占め、台が大きいほど指数的に大きくなる。

\section{共通ランダム測定(CRM)}
\label{sec:crm}

\subsection{定義と不偏性}
CRM\cite{vermersch2024}は、古典計算で作った近似状態 $\sigma$(prior)を使って、実験の影 $\hat\rho^{(r)}$ から「同じ設定 $u^{(r)}$ で $\sigma$ を測ったときの影」$\sigma^{(r)}$ を引き、$\sigma$ を足し戻す:
\begin{equation}
  \hat\rho^{(r)}_\sigma=\hat\rho^{(r)}-\sigma^{(r)}+\sigma .
  \label{eq:crm}
\end{equation}
$\sigma^{(r)}$ は prior を\emph{標本にせず}、古典計算で求めた厳密な出力確率から作る(元論文の式(3))。設定について平均すると $\mathbb E[\hat\rho^{(r)}]=\rho$、$\mathbb E[\sigma^{(r)}]=\sigma$ なので、$\sigma$ をどう選んでも推定は偏らない。

パウリ文字列 $P$ について書けば、1 つの設定での CRM の推定値は
\begin{equation}
  3^{\lvert A\rvert}\,\mathbf 1[\text{一致}]\,\bigl(\bar m-y\bigr)+y,\qquad y\equiv\Tr[P\sigma]
\end{equation}
である($\bar m$ は一致した設定での $\pm1$ の出力の積の平均)。実装に必要なのは prior の期待値 $y$ だけで、$\sigma$ が混合状態でも構わない。

\subsection{分散と利得}
\label{sec:gain}
CRM では、基底のくじの揺らぎが $x$ ではなく外れ $\Delta\equiv x-y$ に比例するようになり、ショット雑音はそのまま残る:
\begin{equation}
  n_u\Var[\text{CRM}]=(3^{\lvert A\rvert}-1)\Delta^2+\frac{3^{\lvert A\rvert}(1-x^2)}{n_m}.
  \label{eq:var-crm}
\end{equation}
これは元論文の単一パウリ文字列の厳密な分散(arXiv 版の付録 B.1 の式(25)、PRX Quantum 版の付録 C.1 の式(C7))である\cite{vermersch2024}。標準の古典シャドウとの分散の比を\emph{利得}と呼ぶ:
\begin{equation}
  \boxed{\ G\equiv\frac{\Var[\text{標準}]}{\Var[\text{CRM}]}=\frac{(3^{\lvert A\rvert}-1)x^2+v_s}{(3^{\lvert A\rvert}-1)\Delta^2+v_s},\qquad v_s=\frac{3^{\lvert A\rvert}(1-x^2)}{n_m}\ }
  \label{eq:gain}
\end{equation}
本論文ではこの比を「利得の式」と呼ぶ。新しい法則ではなく、元論文の式の比である。

利得の式から次のことが読み取れる。
\begin{enumerate}
  \item \textbf{利得を決めるのは観測量ごとの外れ $\Delta$ である。} prior の大域的な忠実度は式に現れない。
  \item \textbf{損得の境目は $\lvert\Delta\rvert=\lvert x\rvert$ である}(元論文 II.C 節)。$y/x$ で言えば $0<y/x<2$ なら得をし、prior の値が真の値の 2 倍を超えるか符号が逆なら損をする。
  \item \textbf{天井がある。} prior が完全($\Delta=0$)でもショット雑音は消えないので
  \begin{equation}
    G_{\max}=1+n_m\frac{(3^{\lvert A\rvert}-1)x^2}{3^{\lvert A\rvert}(1-x^2)}
    \label{eq:gmax}
  \end{equation}
  で頭打ちになる。$\lvert x\rvert\to1$ で天井は発散する。
  \item \textbf{設定の数 $n_u$ は比で約分される。} したがって $G$ は「同じ精度を得るのに必要な設定の数を何分の 1 にできるか」を表す。本論文の数値は特に断らない限り $n_m=100$ での値である。
  \item \textbf{1 設定 1 ショット($n_m=1$)では天井が低い。} 総ショット数を固定すると分散を最小にするのは $n_m=1$ だが、そのときの天井は
  \begin{equation}
    G_{\max}(n_m=1)=\frac{1-x^2/3^{\lvert A\rvert}}{1-x^2}\approx\frac{1}{1-x^2}
  \end{equation}
  にとどまる。CRM の利点は、設定の切り替えにコストがかかり 1 設定あたり多くのショットを取る実験で大きい。元論文も、CRM を有限の設定数による誤差を減らす手法と位置づけている。
\end{enumerate}

相対誤差 $\varepsilon\equiv\Delta/x$ と、基底のくじとショット雑音の比 $g\equiv(3^{\lvert A\rvert}-1)x^2/v_s$ を使うと
\begin{equation}
  G=\frac{g+1}{\varepsilon^2g+1},\qquad\frac1G\simeq\varepsilon^2+\frac1{G_{\max}}
  \label{eq:two-regime}
\end{equation}
と書ける。利得は prior の相対誤差で決まる部分($\varepsilon^2$)と、ショット雑音で決まる部分($1/G_{\max}$)の小さい方に支配される。

\subsection{元論文における忠実度の役割}
元論文は、多コピー観測量(例えば純度やエントロピーの多項式近似)について、Hilbert--Schmidt 距離で書いた分散の上界を与え、純粋な prior $\ket\phi$ なら忠実度 $F=\bra\phi\rho\ket\phi\ge1/2$ を prior の良し悪しの目安に使っている。本研究が扱うのは単一コピーの局所観測量で、観測量の台に縮約した prior は一般に混合状態になる。この設定では忠実度は利得を決めない(第\ref{ch:locality}章)。両者は扱う対象が違うだけで、矛盾しない。

\subsection{和の観測量}
二重占有\eqref{eq:docc-pauli}やスピン相関 $\ev{S^z_iS^z_j}$ のように、複数のパウリ項の和で書かれる観測量では、同じ設定で同時に当たる項の組が分散に寄与する。この場合の厳密な分散は第\ref{ch:methods}章と付録\ref{app:sumvar}で与える。

% =====================================================================
\chapter{手法}
\label{ch:methods}
% =====================================================================

\section{和の観測量の厳密な分散}
\label{sec:sumvar}

単一パウリ文字列の分散\eqref{eq:var-crm}を、パウリ項の和
\begin{equation}
  O=\sum_k c_kP_k,\qquad x_k=\Tr[P_k\rho],\quad y_k=\Tr[P_k\sigma],\quad \Delta_k=x_k-y_k
\end{equation}
に拡張する。2 つの項 $P_k,P_l$ が\emph{両立する}とは、共通の量子ビットの上でパウリ文字が等しいことをいう。両立する項は同じ設定で同時に「当たる」ことができ、その確率は $3^{-\lvert A_k\cup A_l\rvert}$ である。両立しない項は同時には当たらない。

\begin{proposition}\label{prop:sumvar}
1 つの設定と $n_m$ ショットからなる推定量について、
\begin{align}
  n_u\Var[\text{標準}]&=\sum_{k,l\ \text{両立}}c_kc_l\,3^{\lvert A_k\cap A_l\rvert}\Bigl[\frac{\ev{P_kP_l}}{n_m}+\Bigl(1-\frac1{n_m}\Bigr)x_kx_l\Bigr]-\Bigl(\sum_kc_kx_k\Bigr)^2,\label{eq:sumvar-std}\\
  n_u\Var[\text{CRM}]&=\sum_{k,l\ \text{両立}}c_kc_l\,3^{\lvert A_k\cap A_l\rvert}\Bigl[\frac{\ev{P_kP_l}-x_kx_l}{n_m}+\Delta_k\Delta_l\Bigr]-\Bigl(\sum_kc_k\Delta_k\Bigr)^2.\label{eq:sumvar-crm}
\end{align}
ここで $\ev{P_kP_l}=\Tr[P_kP_l\rho]$ で、両立する項の積は共通の量子ビットで恒等になるので、それ自体が台 $A_k\triangle A_l$ のパウリ文字列である。
\end{proposition}
導出は付録\ref{app:sumvar}に示す。単一の項では式\eqref{eq:var-std}・\eqref{eq:var-crm}に戻る。

この式から次のことが分かる。
\begin{itemize}
  \item CRM の分散には、合計の外れ $\sum_kc_k\Delta_k$ ではなく\emph{項ごとの外れ} $\Delta_k\Delta_l$ が入る。合計が同じでも、項ごとの誤差が打ち消し合っている prior は損をする(第\ref{ch:priors}章の二重占有の例)。
  \item 必要なのは真の状態の $\ev{P_kP_l}$ と prior の項ごとの値 $y_k$ だけで、prior の $\ev{P_kP_l}$ は要らない。
  \item ホッピング $\tfrac12(XZ\cdots ZX+YZ\cdots ZY)$ の 2 項は端の量子ビットで文字が違うので両立しない。粒子数を保つ状態では $\ev{XZ\cdots ZX}=\ev{YZ\cdots ZY}$ なので、分散はホッピングの期待値と $\Delta$ だけで書ける。
\end{itemize}
本研究では、この式を Python(\texttt{crm\_pauli\_sum\_variance.py})と Julia(\texttt{crm\_exact\_gain.jl})で実装した。単一パウリ文字列では式\eqref{eq:gain}と機械精度で一致する。和の観測量では、1 次元の $n_m$ 掃引と相図の全 223 点でモンテカルロの模擬と比べ、差が分散の比の標本誤差の範囲に収まることを確かめた(標準化した差の平均 $0.02$ 〜 $0.09$、標準偏差 $0.88$ 〜 $0.99$)。

\section{prior}
\label{sec:priors}

本研究で比べる prior を表\ref{tab:priors}にまとめる。

\begin{table}[t]
\centering\small
\caption{本研究で比べる prior。「厳密な状態」は、作るのに真の基底状態が要るかどうか。}
\label{tab:priors}
\begin{tabular}{@{}llll@{}}
\toprule
prior & 中身 & 厳密な状態 & 対称性 \\
\midrule
UHF & 非制限 Hartree--Fock の Slater 行列式 & 不要 & スピンの向きを選ぶ \\
UHF-sym & UHF をスピン回転で平均した混合状態 & 不要 & 回転不変 \\
PHF & UHF を全スピン 0 に射影した状態 & 不要 & 回転不変 \\
切断 MPS & 厳密な状態を結合次元 $\chi$ に切り詰めたもの & 必要 & 保つことが多い \\
変分 MPS & 結合次元を $\chi$ に制限した DMRG の基底状態 & 不要 & 破ることがある \\
変分 MPS-sym & 変分 MPS をスピン回転で平均した混合状態 & 不要 & 回転不変 \\
\bottomrule
\end{tabular}
\end{table}

\subsection{UHF と Wick の定理}
UHF は、相互作用を平均場 $n_{i\up}n_{i\dn}\to\ev{n_{i\up}}n_{i\dn}+n_{i\up}\ev{n_{i\dn}}-\ev{n_{i\up}}\ev{n_{i\dn}}$ で置き換えた 1 体ハミルトニアンを自己無撞着に解いて得る。初期値を Néel 配置にとると、副格子ごとに磁化を持つ解に落ちる。UHF の状態はスピンごとの Slater 行列式なので、期待値はスピンごとの相関行列 $C^{\sigma}_{pq}=\ev{\cd_{p\sigma}c_{q\sigma}}$ から Wick の定理で求まる。例えば
\begin{align}
  \ev{Z_{p}}&=1-2C_{pp},\\
  \ev{Z_{p}Z_{q}}&=1-2C_{pp}-2C_{qq}+4\bigl(C_{pp}C_{qq}-C_{pq}C_{qp}\bigr)\quad(\text{同じスピン}),\\
  \bigl\langle\tfrac12(XZ\cdots ZX+YZ\cdots ZY)\bigr\rangle&=C_{pq}+C_{qp}
\end{align}
で、行列積状態を経由せずに任意の大きさの系で $O(N^3)$ の計算量で扱える。2 次元の大きな系でも厳密に prior の値を計算できることが、平均場 prior の利点である。

\subsection{スピン回転で平均した prior}
\label{sec:rotavg}
真の基底状態はスピン回転で不変だが、UHF はスピンの向きを 1 つ選んでいる。そこで、全スピンを一様に回して平均した混合状態
\begin{equation}
  \bar\sigma=\int d\Omega\,R(\Omega)\,\sigma\,R(\Omega)^\dagger
\end{equation}
を prior にする。CRM に必要なのは期待値だけなので、混合状態でも構わない。観測量の側を回して平均すればよく、スピン $\Sv_i$ はベクトルとして回るので
\begin{equation}
  \overline{\ev{S^z_i}}=0,\qquad\overline{\ev{S^z_iS^z_j}}=\tfrac13\ev{\Sv_i\cdot\Sv_j}
\end{equation}
となり、$n_i$、$n_in_j$、$n_{i\up}n_{i\dn}$、$\Sv_i\cdot\Sv_j$ は変わらない。$Z_{i\up}=1-n_i-2S^z_i$ を使うと
\begin{equation}
  \ev{Z_{i\up}}_{\bar\sigma}=1-\ev{n_i},\qquad
  \ev{Z_{i\up}Z_{j\up}}_{\bar\sigma}=\bigl\langle(1-n_i)(1-n_j)\bigr\rangle+\tfrac43\ev{\Sv_i\cdot\Sv_j}
  \label{eq:rotavg}
\end{equation}
で、オンサイトの $Z_{i\up}Z_{i\dn}$ は回転で変わらない。UHF に当てはめたものを UHF-sym、変分 MPS に当てはめたものを変分 MPS-sym と呼ぶ。2 次元の計算では、UHF-sym を回転した(非共線の)一般化 Slater 行列式の Wick の定理と極角についての数値積分で求めた。

\subsection{スピン射影 HF(PHF)}
UHF を全スピン $S=0$ の部分空間に射影した状態 $P_{S=0}\ket{\rm UHF}$ である。射影は回転について積分した形 $P_{S=0}\propto\int d\Omega\,R(\Omega)$ に書けるので、期待値は回転した Slater 行列式どうしの重なりの積分(一般化 Wick の定理)で計算できる。回転平均(UHF-sym)が状態の混合であるのに対し、PHF は純粋状態の重ね合わせである。

\subsection{切断 MPS と変分 MPS}
\begin{itemize}
  \item \textbf{切断 MPS}:厳密な(または十分大きな結合次元の)基底状態を、正準形で結合次元 $\chi$ に切り詰めて規格化したもの。$\chi$ が prior の質のつまみになる。厳密な状態が要るので、実用の prior というより「結合次元 $\chi$ で表せるほぼ最良の近似」の目安である。元論文の数値例の prior もこの作り方である。
  \item \textbf{変分 MPS}:結合次元の上限を $\chi$ にした DMRG の基底状態。厳密な状態は要らない。低い $\chi$ の DMRG は局所解に落ちやすいので、Néel の積状態と乱数の MPS の複数の出発点から始め、エネルギーの最も低いものを採った。
\end{itemize}

\section{参照状態}
\label{sec:reference}

量子計算機で用意する状態 $\rho$ の代わりに、ハバード模型の基底状態を古典計算で作り、それを「擬似実験」の状態とした。
\begin{itemize}
  \item \textbf{厳密な参照(80 系)}:1 次元の鎖(開放端 $L=8,10,12,14,16,24,32,48,64$、周期境界 $L=8,\dots,16$)、2 次元の小さな系($2\times4$、$2\times6$ のシリンダーとトーラス、$4\times3$ のシリンダーとトーラス)の 20 通りの格子に、$U=2,4,8,12$ を掛けた 80 系。小さな系は厳密対角化、大きな系は DMRG で作り、エネルギーの分散 $\ev{H^2}-\ev{H}^2$ がすべての系で $10^{-8}$ 以下であることを確かめた。$4\times3$ のトーラスだけは二部格子でない。
  \item \textbf{大きな系}:1 次元 $L=128$(256 量子ビット)、2 次元の $W=4$ と $W=6$ のシリンダー($L_x=8$、64 と 96 量子ビット)は、DMRG の基底状態を使う(結合次元は $L=128$ で $256$、$W=4$ で $192$、$W=6$ で $512$)。$L=128$ の $\chi=256$ は収束していないが、本研究が比べるのは固定した状態 $\rho$ に対する分散の比なので、その状態を「実験で用意された状態」とみなせば議論は閉じる。2 次元の DMRG の状態は完全には収束しておらず、スピン回転の対称性を破っている(第\ref{ch:2d}章で注意する)。
\end{itemize}
DMRG は ITensorMPS\cite{itensor}で行い、大きな計算は研究室の計算機クラスターで実行した。

\section{測定の模擬と利得の評価}
\label{sec:simulation}

利得は 2 通りで評価した。
\begin{enumerate}
  \item \textbf{厳密な式}:単一パウリ文字列は式\eqref{eq:gain}、和の観測量は命題\ref{prop:sumvar}。統計誤差がないので、本論文の表の数値は特に断らない限りこちらである。
  \item \textbf{モンテカルロの模擬}:参照状態の MPS からランダムな基底での測定を実際に標本として引き、標準の推定量と CRM の推定量を同じデータで計算する。$n_u$ 設定・$n_m$ ショットの実験を $n_{\rm rep}$ 回繰り返し、推定値の分散の比を取る。prior 側は標本にせず、各設定での条件付き期待値を厳密に計算する。
\end{enumerate}

\paragraph{窓サンプリング.}観測量の台が連続した量子ビットの窓 $[q_1,q_2]$ に収まるとき、窓の外の量子ビットの測定結果は推定値に現れない。MPS の直交中心を $q_1$ に置くと、窓の左端の Schmidt 指標を基底によらない確率で先に引き、窓の中だけを逐次サンプリングすればよい。これにより、計算量は窓の幅だけで決まり、系の大きさによらない。

\paragraph{モンテカルロの精度についての注意.}分散の比の推定値は、対数で約 $\sqrt{4/(n_{\rm rep}-1)}$ の標準偏差を持つ。$n_{\rm rep}=30$〜$50$ ではこれが $0.28$〜$0.37$ になり、利得が 1.5 倍以上外れることも珍しくない。さらに、同じ測定データで複数の prior を比べると、標準の推定量の分散が全 prior の共通の分子になるので、その標本の偶然で全 prior の値が一緒に膨らむ(または縮む)。本研究の初期の表にはこの影響を受けた値があり、すべて厳密な式の値に置き換えた。

% =====================================================================
\chapter{局所性:大域的な忠実度は利得を決めない}
\label{ch:locality}
% =====================================================================

利得の式\eqref{eq:gain}には prior の大域的な忠実度が現れない。この章では、それが実際の prior でどういう帰結を持つかを確かめ、prior の良さを観測量の台の上で測るにはどの量が適切かを論じる。

\section{系の大きさを変える}
\label{sec:size}

1 次元の鎖($U=4$、半充填)で、系の長さを $L=8,16,32$(16〜64 量子ビット)と変え、鎖の中央の観測量の利得を調べた。prior は参照状態を結合次元 $\chi_p$ に切り詰めた MPS である。表\ref{tab:size}に結果を示す。

\begin{table}[t]
\centering\small
\caption{系の長さと利得(1 次元、$U=4$、$n_m=100$、分散の厳密な式による値)。$F$ は prior の大域的な忠実度。}
\label{tab:size}
\begin{tabular}{@{}ccccc@{}}
\toprule
$L$ & $F$($\chi_p=2$) & 隣接 $Z_\up Z_\up$($\chi_p=2$) & 隣接 $S^zS^z$($\chi_p=2$) & オンサイト $ZZ$($\chi_p=16$) \\
\midrule
8  & 0.35   & 11.1 & 7.5  & 50.9 \\
16 & 0.097  & 13.3 & 10.3 & 50.6 \\
32 & 0.0069 & 14.0 & 12.5 & 50.7 \\
\bottomrule
\end{tabular}
\end{table}

$\chi_p=2$ の prior の大域的な忠実度は $L$ とともに指数的に下がり、$L=32$ では $0.0069$ と、参照状態とほとんど直交している。それでも局所観測量の利得は下がらない。オンサイト $ZZ$ は $\chi_p\ge8$ で天井 $G_{\max}=50.7$ に張り付き、$L$ によらない。

\todo{図:crm\_fig2\_locality.png はモンテカルロの値で描いている。厳密な値で描き直す}

この性質は 256 量子ビット($L=128$)まで、また厳密な参照を持つ 80 系の総当たりでも成り立つ。$U=12$ で $L=8$ から $64$ に広げると、切断 MPS prior の大域的な忠実度は約 2 万分の 1 に下がる一方、観測量の台に縮約した忠実度の変化は $3\times10^{-5}$、利得の変化は 2\% にとどまった。

\paragraph{なぜか.}忠実度は系全体の量で、サイトごとにわずかな誤差 $f<1$ があれば $F\sim f^{N}$ と系の大きさについて指数的に下がる。一方、利得を決める外れ $\Delta=\Tr[P(\rho-\sigma)]$ は観測量の台 $A$ に縮約した密度行列 $\rho_A,\sigma_A$ だけで決まる:
\begin{equation}
  \Delta=\Tr\bigl[P(\rho_A-\sigma_A)\bigr].
\end{equation}
系の大きさが効くのは台の周りの縮約密度行列を通してだけで、それは局所的な性質である。

\section{台に縮約した距離}
\label{sec:tracedist}

\subsection{トレース距離が最良の上界であること}
$\lVert P\rVert_\infty=1$ なので、H\"older の不等式から
\begin{equation}
  \lvert\Delta\rvert\le\lVert\rho_A-\sigma_A\rVert_1=2D_A,\qquad D_A\equiv\tfrac12\lVert\rho_A-\sigma_A\rVert_1
  \label{eq:holder}
\end{equation}
が成り立つ。さらにトレースノルムは作用素ノルムの双対ノルム $\lVert X\rVert_1=\sup_{\lVert O\rVert_\infty\le1}\Tr[OX]$ なので、「$P$ の固有値が $\pm1$ である」という情報だけを使う限り、これより良い上界はない。等号が成り立つのは、$P$ の固有値の符号のパターンが $\rho_A-\sigma_A$ のそれと一致するときである。

\subsection{半充填では等号が成り立つ}
半充填でスピンの $z$ 成分を保つ状態では、サイト $i$ の 2 量子ビットの縮約密度行列は占有数の基底で
\begin{equation}
  \rho_A=\mathrm{diag}\bigl(d,\ \tfrac12-d,\ \tfrac12-d,\ d\bigr)
\end{equation}
と二重占有率 $d$ だけで決まる。差は $\rho_A-\sigma_A=\mathrm{diag}(u,-u,-u,u)$ の形になり、オンサイトの $Z_\up Z_\dn=\mathrm{diag}(+1,-1,-1,+1)$ と符号のパターンが一致する。したがって
\begin{equation}
  \lvert\Delta_{ZZ}\rvert=2D_A
\end{equation}
が厳密に成り立つ。オンサイト $ZZ$ の 588 点で $\lvert\Delta\rvert/(2D_A)$ を調べると、中央値も最大値も $1.0000$ だった。

prior がスピンの $z$ 成分を破ると、差は $\mathrm{diag}(u,-u+s,-u-s,u)$ の 2 つのパラメータを持つ。$r\equiv\lvert\Delta(S^z)\rvert/\lvert\Delta_{ZZ}\rvert=s/(4u)$ とおくと
\begin{equation}
  \frac{\lvert\Delta_{ZZ}\rvert}{2D_A}=\min\Bigl(1,\ \frac{2}{1+4r}\Bigr)
\end{equation}
で、$r\le1/4$ なら等号が保たれる。

\subsection{忠実度ではうまくいかない理由}
台に縮約した prior $\sigma_A$ は、大域的には純粋な prior であっても、台の外とのエンタングルメントのため一般に混合状態になる。このとき忠実度 $F(\rho_A,\sigma_A)=\bigl(\Tr\sqrt{\sqrt{\rho_A}\sigma_A\sqrt{\rho_A}}\bigr)^2$ について次のことが言える。
\begin{itemize}
  \item 差 $\rho_A-\sigma_A$ の関数ではなく、ノルムでもない。
  \item $\rho_A=\sigma_A$ の周りで $1-F$ は差の 2 次から始まる。一方 $\Delta$ は 1 次なので、$1-F$ から $\Delta$ を予測する定数は存在しない(実データで $\lvert\Delta\rvert/(2(1-F_A))$ は 4 桁以上ばらつく)。
  \item Fuchs--van de Graaf の不等式 $1-\sqrt F\le D\le\sqrt{1-F}$ で $D$ を挟んでも、粗い prior($\chi_p=2,4$)では利得の予測が 5〜14 倍の幅を持つ。
\end{itemize}
実際、$L=16$、$U=4$ の鎖のサイト 2 では、UHF の台の忠実度は $0.700$ で $\chi_p=2$ の切断 MPS($0.802$)より低いのに、オンサイト $ZZ$ の利得は $48.6$ 対 $2.25$ と 22 倍良い。UHF の距離はスピンの向きを選ぶこと(上の $s$ の方向)に費やされていて、$ZZ$ の外れには効かないからである。

なお、元論文が忠実度を prior の目安に使うのは、多コピー観測量と純粋な prior の文脈である。純粋な prior では $\lVert\rho-\sigma\rVert_2^2\le2(1-F)$ が成り立ち、分散の上界(Hilbert--Schmidt 距離の 2 乗)と次数も合う。本章の結果はそれと矛盾しない。

\subsection{台に縮約した量の限界}
台の上の量(忠実度でもトレース距離でも)は、台を決めた時点で 1 つの数になる。2 量子ビットの台には恒等でないパウリ文字列が 15 個あり、それぞれ別の利得を持つ。例えば同じサイトの上のオンサイト $ZZ$・二重占有・$n_i$・$S^z_i$ は同じ $D_A$ を共有するが、$\chi_p=4$ の prior で利得はそれぞれ $43.8$、$26.9$、$1.00$、$1.00$ である。トレース距離が与えるのは台の上の全観測量に共通の最悪値であり、観測量ごとの利得を決めるのは結局 $\Delta$ そのものである。

\section{まとめ}
\begin{itemize}
  \item prior の大域的な忠実度は局所観測量の利得を決めない。系を 256 量子ビットまで広げても利得は保たれる。
  \item 観測量の台の上で prior の良さを測るなら、トレース距離が $\lvert\Delta\rvert$ の最良の上界で、半充填のオンサイト量では等号が成り立つ。
  \item ただし観測量ごとの利得は $\Delta$ で決まる。次章では、ハバード模型の現実的な prior について $\Delta$ が観測量ごとにどう振る舞うかを調べる。
\end{itemize}

% =====================================================================
\chapter{観測量ごとの外れと prior の物理}
\label{ch:priors}
% =====================================================================

前章で、利得を決めるのは観測量ごとの外れ $\Delta$ であることを見た。この章では、ハバード模型の現実的な prior について $\Delta$ が観測量ごとにどう振る舞うか、それがなぜかを調べる。データは厳密な参照を持つ 80 系(第\ref{sec:reference}節)の中央のサイトとその隣で、利得は $n_m=100$ での厳密な値である。

\section{全体像}

表\ref{tab:overview}に、$U=12$ の 20 系での利得の中央値をまとめる。

\begin{table}[t]
\centering\small
\caption{$U=12$、厳密な参照を持つ 20 系での利得の中央値($n_m=100$)。二重占有は 4 項の和として測るときの厳密な値。「天井」は prior が完全な場合の値。}
\label{tab:overview}
\begin{tabular}{@{}lccccccc@{}}
\toprule
観測量 & UHF & UHF-sym & PHF & 切断 $\chi{=}4$ & 変分 $\chi{=}8$ & 変分 $\chi{=}8$-sym & 天井 \\
\midrule
オンサイト $ZZ$(電荷) & 465 & 465 & 513 & 151 & 555 & 555 & 555 \\
隣接 $Z_\up Z_\up$(スピン) & 1.2 & 7.0 & 14.2 & 29.2 & 16.8 & 19.8 & 31.6 \\
単一サイト $Z_\up$ & 0.02 & 1.00 & 1.00 & 1.00 & 0.11 & 1.00 & 1 \\
二重占有(4 項の和) & 1.8 & 123 & 127 & 78 & 13.3 & 129 & --- \\
\bottomrule
\end{tabular}
\end{table}

この表から 3 つのことが読み取れる。
\begin{enumerate}
  \item \textbf{観測量によって prior の得意・不得意が逆転する。} UHF は電荷の量(オンサイト $ZZ$)では天井の 8 割に届くが、スピンの量では損をする。切断 MPS($\chi=4$)はその逆である。
  \item \textbf{対称性を回復すると損が消える。} UHF も変分 MPS も、スピン回転で平均すると単一サイト $Z_\up$ と二重占有の損が消える。
  \item \textbf{電荷とスピンの両方で良いのは、変分 MPS を回転平均したものである。}
\end{enumerate}
以下、それぞれの理由を説明する。

\section{電荷の量:強結合で平均場の誤差が消える}
\label{sec:charge}

\subsection{二重占有の誤差は 28\% で止まる}
強結合での真の二重占有率は、式\eqref{eq:d-strong}より 1 次元で
\begin{equation}
  d_{\rm 真}\approx\frac{4t^2}{U^2}\Bigl(\frac14-\ev{\Sv_i\cdot\Sv_j}\Bigr)
\end{equation}
で、隣の結合が一重項を組む確率に比例する($U=12$ で予測 $0.0188$、実測 $0.0180$)。一方、共線 UHF は $\up$ と $\dn$ が独立な Slater 行列式の積なので
\begin{equation}
  d_{\rm UHF}=\ev{n_{i\up}}\ev{n_{i\dn}}=\frac14-m^2
\end{equation}
($m$ は副格子磁化)である。自己無撞着条件を $U\gg t$ で展開すると $m\approx\frac12-\frac{2t^2}{U^2}$ となり
\begin{equation}
  d_{\rm UHF}\approx\frac{2t^2}{U^2}=\frac{4t^2}{U^2}\Bigl(\frac14+\frac14\Bigr).
\end{equation}
括弧の中は古典的な反平行の積状態 $\ket{\up\dn}$ の一重項成分 $\frac14-(-\frac14)=\frac12$、つまり「半分だけ一重項」である。真の値は無限の Heisenberg 鎖の $\ev{\Sv_i\cdot\Sv_j}=\frac14-\ln2$ を使って $\frac14+(\ln2-\frac14)=0.693$ なので、二重占有の相対誤差は
\begin{equation}
  \delta\equiv\frac{d_{\rm UHF}}{d_{\rm 真}}-1\longrightarrow\frac{0.5}{0.693}-1=-0.28
\end{equation}
と、$U$ によらない定数に止まる。

\subsection{オンサイト $ZZ$ に対する誤差は $(t/U)^2$ で消える}
ところが利得を決めるのは、オンサイト $ZZ$ $=4d-1$ に対する相対誤差である:
\begin{equation}
  \varepsilon_{ZZ}=\frac{4\lvert d_{\rm 真}-d_{\rm UHF}\rvert}{\lvert4d_{\rm 真}-1\rvert}\approx4d_{\rm 真}\lvert\delta\rvert\approx3.1\Bigl(\frac tU\Bigr)^2.
\end{equation}
分子は $d\propto(t/U)^2$ とともに 0 へ向かい、分母は 1 に張り付く。\textbf{二重占有そのものの誤差は 28\% のまま残るのに、オンサイト $ZZ$ に対する誤差は $(t/U)^2$ で消える}。これが「平均場 prior は電荷の量では強結合で良くなる」ことの正体である。同時に、天井 $G_{\max}\propto x^2/(1-x^2)$ は $1-x^2\approx8d\propto(t/U)^2$ で上がるので、オンサイト $ZZ$ の利得は強結合で大きくなる。

\subsection{弱結合側}
弱結合では、真の状態は相関で二重占有を抑えるのに、UHF は $\up$ と $\dn$ が無相関なので二重占有を過大評価する。$U\simeq4$ で過大評価と過小評価が入れ替わり、誤差がいったん小さくなる。

\subsection{電荷の量が境界条件に鈍い理由}
半充填の二部格子では、粒子正孔対称性によりサイトごとの電子数が厳密に $\ev{n_i}=1$ になる(付録\ref{app:ph})。開放端の影響は電荷の密度には現れず、結合の強弱の交替として現れる。磁化を持つ UHF も、粒子正孔変換とスピン回転を組み合わせた変換の対称性を保つので $\ev{n_i}=1$ を満たす。電荷の量で UHF が境界条件や系の大きさによらず良い prior であるのはこのためである。二部格子でない $4\times3$ トーラスでは UHF の電荷が 6.5\% 偏っていた。

\section{スピンの量:平均場が損をする 2 つの理由}
\label{sec:spin}

\subsection{理由 1:スピンの向きを選ぶ}
有限系の真の基底状態はスピン回転で不変で、$\ev{S^z_i}=0$ である。共線 UHF は副格子ごとに $\ev{S^z_i}=\pm m$ を持つので、単一サイトの $Z_{i\up}=1-n_i-2S^z_i$ を大きく外す。真の値が 0 なので損得の境目を必ず越え、$U=12$ では 20 系すべてで $G=0.02$ 程度の大損になる。隣接 $Z_\up Z_\up$ でも、スピンの向きを選ぶことで $z$ 方向の相関を過大評価し、80 系中 17 系で損をする。

この原因は、スピン回転で平均した UHF-sym(第\ref{sec:rotavg}節)で取り除ける。単一サイト $Z_\up$ は $1-\ev{n_i}$ になって外れが消え、隣接 $Z_\up Z_\up$ で損をする系は 17 から 0 になる。

\subsection{理由 2:量子的な一重項の相関が足りない}
回転平均は $\ev{\Sv_i\cdot\Sv_j}$ を変えない。UHF の $\ev{\Sv_i\cdot\Sv_j}$ は古典的な Néel 配置の値 $-1/4$ で、真の値($1$ 次元の無限鎖で $\frac14-\ln2=-0.443$)の量子的な一重項の相関を持たない。隣接 $\ev{S^z_iS^z_j}$ について、真の値は $\ev{\Sv_i\cdot\Sv_j}/3$、共線 UHF は $-m^2\to-1/4$ なので
\begin{equation}
  \frac yx\longrightarrow\frac{3/4}{\lvert\ev{\Sv_i\cdot\Sv_j}\rvert},\qquad
  \varepsilon\longrightarrow\Bigl\lvert1-\frac{3/4}{\lvert\ev{\Sv_i\cdot\Sv_j}\rvert}\Bigr\rvert
\end{equation}
となる。1 次元では $\varepsilon\to0.69$ で、実測($L=12$ の周期境界)の $0.690$($U=8$)、$0.684$($U=12$)と一致する。損得の境目 $y/x=2$ は $\lvert\ev{\Sv_i\cdot\Sv_j}\rvert=3/8$ に当たり、1 次元(0.443)では得、2 次元の正方格子(約 0.335)では損になる。

この原因は回転平均では直らない。UHF-sym の隣接 $Z_\up Z_\up$ の利得は $U=12$ で 7.0 と、天井 31.6 に遠く届かない。

\section{スピン射影 HF:効果は $1/L$ で消える}
\label{sec:phf}

量子的な相関を足す方法として、UHF を全スピン 0 に射影した PHF を試した。表\ref{tab:overview}のとおり、PHF の隣接 $Z_\up Z_\up$ は 14.2 と UHF-sym の 2 倍だが、天井には届かない。しかも\textbf{この改善は系の大きさとともに消える}。

射影の効果は回転子の描像で説明できる。UHF の 2 つの副格子のスピンはそれぞれ大きさ $\sim mL/2$ の大きなスピンにまとまっている。射影はそれらを全スピン 0 に組み直すが、1 つの結合に割り当てられる量子的な相関は $1/L$ に比例する。実際、射影による隣の $\ev{\Sv_i\cdot\Sv_j}$ の変化は
\begin{equation}
  L\,\Delta\!\ev{\Sv_i\cdot\Sv_j}=-2m
\end{equation}
で予測でき($U=12$、$L=64$ で実測 $-0.963$ 対予測 $-0.972$)、エネルギーの変化も $E_{\rm PHF}-E_{\rm UHF}=-2mJ$ と $L$ によらない。$L=64$ では、UHF と真の値の差のうち射影で埋まるのは 11\% だけである。数十〜数百量子ビットの系で PHF を使う利点はない。

\section{行列積状態の prior:作り方で性質が逆になる}
\label{sec:mps}

低い結合次元 $\chi$ の MPS は、結合ごとに量子的な相関を持てる。ところが、切断 MPS と変分 MPS では性質が逆になった(表\ref{tab:overview})。
\begin{itemize}
  \item \textbf{切断 MPS($\chi=4$)}:隣接 $Z_\up Z_\up$ の利得は 29.2 と天井の 9 割に届くが、二重占有は 78、オンサイト $ZZ$ は 151 にとどまる。
  \item \textbf{変分 MPS($\chi=4$)}:オンサイト $ZZ$ は 553 と天井に届き、相関エネルギーの 78〜88\% を $L$ によらず取り戻すが、中央サイトの $\lvert\ev{S^z}\rvert$ が約 0.30 あり、単一サイト $Z_\up$ では 20 系すべてで損をする。
\end{itemize}
中央の結合で量子数ごとに Schmidt 分解すると、この違いは少ない状態の使い方の違いだと分かる。厳密な状態($L=64$、$U=12$)では電荷がゆらいだ成分の重みが 1.7\% ある。切断 MPS は重みの大きいスピンの成分だけを上下対称に残し、この電荷の成分を捨てる。変分 MPS は電荷の成分を残してその重み(1.8\%)を正しく再現する代わりに、スピンの成分を一方の向きだけにする。

\section{変分 MPS の回転平均:電荷とスピンの両方で良い prior}
\label{sec:varsym}

変分 MPS の弱点はスピンの向きを選ぶことだけなので、UHF-sym と同じ回転平均(式\eqref{eq:rotavg})で直せるはずである。clara に保存した変分 MPS(80 系 × 4 つの $U$ × $\chi=2,4,8,16$ の 320 個)について、回転平均に要る $\ev{n_i}$、$\ev{n_in_j}$、$\ev{\Sv_i\cdot\Sv_j}$ を求めた。

結果を表\ref{tab:varsym}に示す。
\begin{table}[t]
\centering\small
\caption{変分 MPS とその回転平均(-sym)の利得の中央値($n_m=100$、各 $U$ の 20 系)。括弧は損をする系の数。}
\label{tab:varsym}
\begin{tabular}{@{}llcccc@{}}
\toprule
$U$ & 観測量 & UHF-sym & 変分 $\chi{=}8$ & 変分 $\chi{=}8$-sym & 変分 $\chi{=}16$-sym \\
\midrule
4  & 隣接 $Z_\up Z_\up$ & 11.2 & 9.71 (1) & 11.7 & 18.9 \\
4  & 二重占有 & 28.0 & 8.67 (3) & 28.0 & 28.1 \\
12 & 隣接 $Z_\up Z_\up$ & 6.95 & 16.8 (1) & 19.8 & 26.7 \\
12 & 二重占有 & 123 & 13.3 & 129 & 129 \\
\bottomrule
\end{tabular}
\end{table}

\begin{itemize}
  \item $\chi\ge8$ の回転平均は、80 系 × 4 観測量(オンサイト $ZZ$、隣接 $Z_\up Z_\up$、単一サイト $Z_\up$、二重占有)のどれでも損をしなかった(単一サイト $Z_\up$ の 2 系は UHF-sym でも損をする観測量の限界)。
  \item 隣接 $Z_\up Z_\up$ では $U\ge8$ で UHF-sym の 2〜3 倍得をする。回転平均した prior の質は $\ev{\Sv_i\cdot\Sv_j}$ の正確さで決まり、その相対誤差は $U=12$ で UHF 0.35 に対して変分 $\chi=8$ で 0.11、$\chi=16$ で 0.05 だった。
  \item 必要なのは変分 MPS を 1 回作ることと期待値の計算だけで、厳密な状態は要らない。厳密な状態を使わない prior の中では、本研究で試した中で最も良い。
\end{itemize}

\section{和の観測量:項ごとの外れ}
\label{sec:sumobs}

二重占有を 4 項の和\eqref{eq:docc-pauli}として測ると、共線 UHF の利得は $U=12$ で 1.8 しかない(表\ref{tab:overview})。二重占有の合計の外れは小さい(相対 28\%、かつ $d$ 自体が小さい)のに、である。命題\ref{prop:sumvar}によれば、分散には項ごとの外れ $\Delta_k\Delta_l$ が入る。共線 UHF の $\ev{Z_{i\up}}=\pm2m$ は真の値 0 から大きく外れ、$Z_{i\up}$ と $Z_{i\dn}$ の外れは符号が逆で合計では打ち消し合うが、分散には打ち消されずに残る。回転平均すると単一の $Z$ の外れが消え、利得は 123 になる。

\textbf{和の観測量では、合計の値が正しいだけでは足りない。各項の値が正しい必要がある。}逆に、参照状態の側が対称性を破っているとき(2 次元の DMRG の状態など)は、回転平均した prior の方が項ごとの外れが大きくなり、利得が下がることがある(第\ref{ch:2d}章)。

\section{まとめ}
\begin{itemize}
  \item 平均場 prior は、電荷の量(オンサイト $ZZ$)では強結合で誤差が $(t/U)^2$ で消え、天井近くまで得をする。二重占有そのものの誤差は 28\% で止まるが、オンサイト $ZZ$ では分母が 1 に張り付くので効かない。
  \item スピンの量で平均場が損をする原因は 2 つある。スピンの向きを選ぶこと(回転平均で直る)と、量子的な一重項の相関が足りないこと(直らない。2 次元では損になる)。
  \item スピン射影は量子的な相関を足すが、その効果は $1/L$ で消える。
  \item 変分 MPS を回転平均すると、電荷とスピンの両方で損をしない prior になる。
  \item 和の観測量では、項ごとの外れが利得を決める。
\end{itemize}

% =====================================================================
\chapter{2 次元とドープ系}
\label{ch:2d}
% =====================================================================

1 次元では、低い結合次元の MPS が局所観測量に十分な prior になる。2 次元の幅 $W$ のシリンダーでは、MPS で状態を表すのに必要な結合次元が幅について指数的に増える($\chi\sim e^{cW}$)。この章では、MPS prior が高くつく 2 次元で、MPS を使わない平均場 prior がどこまで代わりになるかを調べる。

\section{設定}
\begin{itemize}
  \item 系:$W=4$ と $W=6$ のシリンダー($L_x=8$、円周方向は周期、軸方向は開放)。64 量子ビットと 96 量子ビット。$U=4,8$。
  \item 参照状態:DMRG の基底状態(結合次元 \todo{W=4: 192、W=6: 512 を確認})。2 次元の DMRG の状態は完全には収束しておらず、副格子の磁化を持つ(スピン回転の対称性を破る)。
  \item prior:UHF、UHF-sym、参照状態を $\chi$ に切り詰めた MPS。
  \item 観測量:中央のサイトの周りのオンサイト $ZZ$、隣接と距離 2 の $Z_\up Z_\up$、二重占有、$S^zS^z$、$y$ 方向のホッピング。
  \item 利得は分散の厳密な式による($n_m=100$)。
\end{itemize}

\section{$W=4$ シリンダー}
\label{sec:w4}

\begin{table}[t]
\centering\small
\caption{$W=4$ シリンダー、$U=4$ の利得(分散の厳密な式、$n_m=100$)。}
\label{tab:w4}
\begin{tabular}{@{}lcccc@{}}
\toprule
観測量 & UHF & UHF-sym & 切断 $\chi{=}16$ & 切断 $\chi{=}64$ \\
\midrule
オンサイト $ZZ$ & 27.0 & 27.0 & 13.0 & 28.6 \\
二重占有 & 6.9 & 3.4 & 9.7 & 23.8 \\
ホッピング($y$) & 12.9 & 12.9 & 16.5 & 16.5 \\
$Z_\up Z_\up$($x$ 方向の隣) & 0.97 & 8.68 & 6.8 & 8.9 \\
$S^zS^z$($x$ 方向の隣) & 0.88 & 6.59 & 4.5 & 7.8 \\
\bottomrule
\end{tabular}
\end{table}

表\ref{tab:w4}から次のことが分かる。
\begin{itemize}
  \item \textbf{オンサイト $ZZ$ とホッピングでは、UHF が $\chi=64$ の MPS に並ぶ。} $U=8$ ではオンサイト $ZZ$ で UHF 141 対 $\chi=64$ 143(天井 145)である。2 次元では MPS の忠実度が急に下がる($U=4$、$\chi=8$ で $F=0.055$。1 次元 $L=16$ では同じ $\chi$ で 0.95)ので、計算量 $O(N^3)$ の平均場がこれに並ぶことには実用上の意味がある。
  \item \textbf{スピン相関では UHF は損をし、UHF-sym で $\chi=64$ に近い水準まで直る。}
  \item \textbf{二重占有では UHF-sym の方が悪い}(6.9 → 3.4、$U=8$ では 10.5 → 5.1)。第\ref{sec:sumobs}節のとおり利得は項ごとの外れで決まる。参照状態の DMRG がスピン回転の対称性を破って $\ev{Z_{i\up}}\ne0$ を持つので、項ごとには共線の UHF の方が合う。これは参照状態の性質によるもので、実験の状態が対称性を保つなら逆になるはずである。
\end{itemize}

\paragraph{モンテカルロの値との違い.}以前の表は反復 50 回のモンテカルロの値で、オンサイト $ZZ$ が 44.0(UHF)・44.5($\chi=64$)と、天井 28.9 を大きく超えていた。同じ測定データで全 prior を評価したため、標準の推定量の分散の偶然の大きさが全列に共通にかかった。同じ乱数の種で再実行すると同じ外れ方を再現したので、計算の誤りではない。

\section{$W=6$ シリンダー}
\label{sec:w6}
\todo{crm\_2d\_w6\_results\_U4p0.tsv / U8p0.tsv の G\_exact 列で表を作る。以前の要約の「UHF $G=157$ 対 $\chi=32$ の $G=96$」(モンテカルロ、反復 30 回)を置き換える}

\section{ドープ系}
\label{sec:doped}
\todo{crm\_2d\_doped\_results\_U4p0/U8p0.tsv(W=4、ホール 4 個)と crm\_2d\_w6\_doped\_results\_*.tsv(W=6、ホール 6 個、窓の位置 x0=0, 2)の G\_exact で表を作る}

$1/8$ のホールを入れると、平均場は電荷の川(ストライプ)を自発的に作るが、その電荷の変調が真の状態より鋭すぎる。サイトごとの密度 $\ev{n_i}$ では UHF の外れが $0.13$〜$0.16$ に達し、平均場 prior はこの観測量では役に立たない。\todo{数値は再計算の結果で確認する}

\section{まとめ}
\todo{再計算の結果がそろってから書く}

% =====================================================================
\chapter{実用上の問題}
\label{ch:practical}
% =====================================================================

この章では、CRM を実際の実験に使うときに問題になる点を扱う。ショットの配分、prior が悪いときに損をしない推定量、読み出し誤差、台の長い観測量(matchgate シャドウ)、他の測定法との比較である。

\section{ショットの配分}
\label{sec:allocation}

総ショット数 $B=n_un_m$ を固定すると、CRM の分散は
\begin{equation}
  \Var[\text{CRM}]=\frac1B\Bigl[(3^{\lvert A\rvert}-1)\Delta^2n_m+3^{\lvert A\rvert}(1-x^2)\Bigr]
\end{equation}
で、$n_m$ について単調に増える(標準のシャドウも同じ)。すなわち\textbf{分散を最小にするのは $n_m=1$(設定を最大限多く)}である。一方、利得 $G$ は $n_m$ とともに伸びる。$G$ は分散の比なので、$G$ を最大にする配分と分散を最小にする配分は違う。

両者の関係は、ショット雑音が支配的になる境目
\begin{equation}
  n_m^\ast=\frac{3^{\lvert A\rvert}(1-x^2)}{(3^{\lvert A\rvert}-1)\Delta^2},\qquad\frac{G}{G_{\max}}=\frac{1}{1+n_m/n_m^\ast}
\end{equation}
で整理できる。$n_m\lesssim n_m^\ast$ なら、$n_m$ を増やしても CRM の分散はほとんど悪化せず、$G$ だけが伸びる。つまり設定の数を減らしてよい。$n_m\gg n_m^\ast$ では分散が線形に悪化する。

例えば $\chi_p=2$ の粗い prior では、$n_m=100$ は最適の 20 倍以上大きく、分散で 10.9 倍(標準偏差で 3.3 倍)損をしている。\textbf{粗い prior ほど $n_m$ を小さくすべきである}。CRM の価値が最も高いのは、設定の切り替えにコストがかかり、大きな $n_m$ を取らざるを得ない実験である。

\section{損をしない推定量:制御変量法としての一般化}
\label{sec:optbeta}

CRM を係数 $\beta$ 付きの制御変量法
\begin{equation}
  \hat o(\beta)=\bar m_\rho-\beta\bigl(\bar Y-\Tr[O\sigma]\bigr),\qquad Y(u)=\mathbb E[X_\sigma\mid u]
\end{equation}
と見る。$\mathbb E[\bar Y]=\Tr[O\sigma]$ なので任意の $\beta$ で不偏であり、通常の CRM は $\beta=1$ の場合である。分散は $\beta$ の 2 次式で、最適な $\beta^\ast=\Cov(\bar m_\rho,\bar Y)/\Var(\bar Y)$ では
\begin{equation}
  G_{\rm opt}=\frac{1}{1-\mathrm{corr}^2(\bar m_\rho,\bar Y)}\ge1
\end{equation}
となり、prior がどれほど悪くても損をしない。

実際には $\beta^\ast$ を同じデータから推定する必要があり、次の問題が出る。
\begin{itemize}
  \item \textbf{$\Var(Y)$ が 0 に近いと推定が発散する。} 例えば $\chi_p=2$ の prior は二重占有を厳密に 0 と予言するので $Y\equiv0$ となる。$\Var(Y)$ は prior の期待値だけから閉じた形で計算できるので、標本で推定する必要はない。
  \item \textbf{推定した係数は推定量に偏りを入れる。} 分散の比だけで比べると、偏りのある推定量を不当に有利に評価する。平均二乗誤差で比べる必要がある。
  \item \textbf{単一のパウリ文字列では prior が約分される。} $Y(u)=\mathbf 1[\text{一致}]\,3^{\lvert A\rvert}y$ なので、最適係数の推定量は prior によらず、「基底が一致したかどうか」という既知の指標だけを制御変量に使う。
  \item \textbf{良い prior では $\beta=1$ の方が良い。} 係数を推定する誤差がなくなる分だけ有利である。
\end{itemize}
$W=4$、$U=8$ のシリンダー(反復 400 回のモンテカルロ)では、$\beta=1$ で損をしていた観測量がすべて直った(UHF prior の $Z_\up Z_\up$ で、$y$ 方向の距離 2 が $0.70\to10.7$、$x$ 方向の隣が $1.35\to17.1$)。一方、prior の良いオンサイト $ZZ$ では $\beta=1$ の方が良かった(分散の比で $169$ 対 $117$)。UHF と UHF-sym の 2 つの prior を回帰で組み合わせると、二重占有で $G=82$ に達し、$\chi=32$ の MPS(82)と並んだ。MPS を作らずに MPS と同じ利得が得られる。

\section{読み出し誤差}
\label{sec:readout}

各量子ビットの読み出しが独立に確率 $p$ で反転すると、台 $\lvert A\rvert$ のパウリ文字列の 1 ショットの平均は $f\equiv(1-2p)^{\lvert A\rvert}$ 倍に縮む。比べる推定量は、どれも同じ量 $fx$ を推定する(最後に $f$ で割る補正は共通なので、分散の比は変わらない):
\begin{itemize}
  \item 標準:$3^{\lvert A\rvert}\bar m$
  \item CRM(そのまま):$3^{\lvert A\rvert}(\bar m-y)+y$ — 誤差のないシミュレーションで作った prior の値を引く。
  \item CRM(誤差込み):$3^{\lvert A\rvert}(\bar m-fy)+fy$ — prior にも同じ誤差モデルを掛ける。
\end{itemize}
利得は式\eqref{eq:gain}で $x\to fx$ とし、外れを $\Delta'=fx-y$(そのまま)または $f(x-y)$(誤差込み)に置き換えたものになる。ビット反転を直接まねたモンテカルロと 1〜2\% で一致することを確かめた。

\begin{table}[t]
\centering\small
\caption{読み出し誤差と利得(1 次元 $L=64$、$U=12$、UHF prior、$n_m=100$)。各セルは「そのまま / 誤差込み」。$P_\ell$ はサイト 17 から $\ell$ サイトのブロックの電荷の偶奇(台 $2\ell$)。}
\label{tab:readout}
\begin{tabular}{@{}lcccccc@{}}
\toprule
観測量 & $\lvert A\rvert$ & $p=0$ & $p=0.002$ & $p=0.005$ & $p=0.01$ & $p=0.02$ \\
\midrule
オンサイト $ZZ$ & 2 & 472 & 372 / 430 & 262 / 377 & 154 / 312 & 67 / 226 \\
$P_4$ & 8 & 597 & 225 / 411 & 76 / 273 & 23.6 / 166 & 5.7 / 82.5 \\
$P_{16}$ & 32 & 591 & 35.6 / 199 & 5.9 / 83.9 & 1.1 / 32.2 & 0.1 / 7.8 \\
\bottomrule
\end{tabular}
\end{table}

結果を表\ref{tab:readout}に示す。
\begin{itemize}
  \item \textbf{そのままの prior は、台の大きい観測量で損に転じる。} 損得の境目 $\lvert fx-y\rvert=\lvert fx\rvert$ は、prior がほぼ正しいとき $f=1/2$、すなわち
  \begin{equation}
    p^\ast\approx\frac{\ln2}{2\lvert A\rvert}\approx\frac{0.35}{\lvert A\rvert}
  \end{equation}
  で、台 32 では $p\approx1\%$ で利得が消える。
  \item \textbf{誤差込みの prior なら損はしないが、天井が下がる。} 測った値が $fx$ に縮むので、基底のくじの揺らぎ $\propto f^2x^2$ が小さくなり、ショット雑音 $\propto1-f^2x^2$ が大きくなる。
  \item 誤差込みの prior には誤差率の推定値 $\hat p$ が要り、台 32 では $\hat p$ が 20\% ずれると利得が 3〜4 割落ちる。必要な較正の精度は台の大きさに比例する。
\end{itemize}
扱ったのは対称で独立なビット反転だけで、非対称な読み出し誤差(混同行列)や状態の用意の誤差は今後の課題である。

\section{台の長い観測量と matchgate シャドウ}
\label{sec:matchgate}

2 次元のシリンダーを蛇行の順に並べると、$x$ 方向のホッピングは JW 弦のため台が $2W+1$ になる($W=4$ で 9)。基底が台の上で全部一致する確率は $3^{-9}\approx1/2$ 万で、50 設定(5000 ショット)の間に一度もサンプルされなかった。CRM を付けても prior の値をそのまま返すだけである。

matchgate(フェルミオンの Gauss 変換)シャドウ\cite{zhao2021,wan2023}は、ランダムな $Q\in SO(2n)$ で表されるフェルミオンの線形変換を掛けてから占有数を測る。ホッピングは Majorana 演算子の 2 次式なので、空間的な距離によらず同じ精度で測れる。prior が Slater 行列式なら、CRM の prior 側の値は Pfaffian で厳密に計算できる。$W=4$、$L_x=2$ のシリンダー(16 量子ビット)で、全エネルギーの推定誤差は標準のパウリシャドウの 11.5 に対し matchgate では 0.87 と約 13 分の 1 になった。

一方、matchgate シャドウに CRM を加えても利得は小さかった($n_m$ を 3000 まで増やしても $G\lesssim1.9$)。Gauss 変換は全モードを混ぜるので、どの設定もすべての 2 次の観測量を少しずつ測り、設定の間の揺らぎ(CRM が打ち消す量)が最初から小さいためである。\textbf{CRM の利得は測定のアンサンブルに依存する。} JW 弦の問題は matchgate で、ランダムなパウリ測定の設定の間の揺らぎは CRM で、と別々の問題に対する道具である。

\section{固定した観測量の組に対する比較}
\label{sec:derand}

あらかじめ決まった観測量の組を測るだけなら、ランダムに基底を選ぶ必要はなく、観測量に合わせて基底を選ぶ方法(derandomization\cite{huang2021})の方が強い。1 次元 $L=16$、$U=4$ で同じ総ショット数で比べると、構造のある 23 個の観測量では誤差の中央値/最悪値が、標準のシャドウ $0.059/0.351$、シャドウ + CRM $0.012/0.054$、derandomization $0.0085/0.028$ だった。ランダムに選んだ 150 個の観測量では、CRM の中央値の改善は消えるが最悪値は半分になる($0.208\to0.096$)。期待値がほぼ 0 の観測量には CRM で減らせる揺らぎがもともとなく、誤差を支配する期待値の大きい少数の観測量で CRM が効くためである。ランダムな測定の利点(測定の後で観測量を選べること)が必要な場面で、CRM はその弱点を補う。

\section{まとめ}
\begin{itemize}
  \item 利得を最大にする配分と分散を最小にする配分は違う。粗い prior ほど 1 設定あたりのショットを減らすべきである。
  \item 制御変量法として係数を最適化すると原理的には損をしないが、推定の誤差と偏りに注意が要る。
  \item 読み出し誤差は台の大きい観測量で CRM を標準より悪くしうる。誤差モデルを prior に掛ける必要がある。
  \item 台の長い観測量には matchgate シャドウが要るが、そこでは CRM の利得は小さい。
\end{itemize}

% =====================================================================
\chapter{結論と今後の課題}
\label{ch:conclusion}
% =====================================================================

\section{結論}

共通ランダム測定(CRM)をフェルミオンのハバード模型に適用し、どんな prior がどの観測量でどれだけ効くか、それはなぜかを調べた。

\begin{enumerate}
  \item \textbf{利得を決めるのは観測量ごとの外れである。} これは元論文の分散の式から出る性質で、本研究はそれがハバード模型の prior でどういう帰結を持つかを確かめた。prior の大域的な忠実度は系の大きさとともに崩れるが、局所観測量の利得は 256 量子ビットまで保たれる。台に縮約した量で prior の良さを測るならトレース距離が最良で、半充填のオンサイト量では $\lvert\Delta\rvert=2D_A$ が厳密に成り立つ。
  \item \textbf{平均場 prior は電荷の量に強く、スピンの量に弱い。} オンサイトの $ZZ$ 相関に対する誤差は強結合で $3.1(t/U)^2$ と消え、利得は天井の 8 割以上に届く。スピンの量では、スピンの向きを選ぶことと量子的な一重項の相関が足りないことの 2 つの理由で損をする。前者はスピン回転の平均で直り、後者は直らない(損得の境目は $\lvert\ev{\Sv_i\cdot\Sv_j}\rvert=3/8$ で、2 次元の正方格子では損になる)。
  \item \textbf{量子的な相関を足す方法の比較。} スピン射影の効果は $1/L$ で消える(回転子の描像で定量的に説明できる)。結合次元を制限した変分 MPS は電荷の量で最良だがスピンの向きを選ぶ。これを回転平均すると、電荷とスピンの両方で損をしない prior になり、厳密な状態を使わない prior の中で最も良い。
  \item \textbf{和の観測量では項ごとの外れが効く。} 合計の値が正しいだけでは足りず、各パウリ項の値が正しい必要がある。和の観測量の厳密な分散を書き下し、実装した。
  \item \textbf{2 次元では平均場 prior が MPS prior に代われる場面がある。} オンサイトの量とホッピングでは、$W=4$ のシリンダーで UHF が $\chi=64$ の MPS に並ぶ。\todo{W=6 とドープ系の結論を再計算の結果で書く}
  \item \textbf{実用上の注意。} 利得を最大にするショットの配分と分散を最小にする配分は違う。読み出し誤差は台の大きい観測量で CRM を標準より悪くしうるので、誤差モデルを prior に掛ける必要がある。
\end{enumerate}

\section{今後の課題}
\begin{itemize}
  \item \textbf{実機の誤差に強い prior。} 非対称な読み出し誤差(混同行列)や、状態の用意の誤差(ゲートの雑音)を prior の側に取り込む方法と、その効果の評価。
  \item \textbf{基底の選び方を prior で偏らせる方法との組み合わせ。} 1 設定 1 ショットでの利得の天井 $\approx1/(1-x^2)$ を超えられるか。
  \item \textbf{データから prior を作る。} 測定データを 2 つに分け、一方で prior を合わせ込み、他方で偏りなく推定する方法。
  \item \textbf{2 次元のより大きな系とドープ系。} 平均場が壊れる領域で、どの prior が有効か。
  \item \textbf{実機での実証。}
\end{itemize}


\appendix
% =====================================================================
\chapter{半充填でサイトごとの電子数が 1 になる理由}
\label{app:ph}
% =====================================================================

開放端の系では一般に、端から電子の密度の振動(Friedel 振動)が立つ。それでも半充填のハバード模型の基底状態では、端でも中央でも厳密に $\ev{n_i}=1$ になる。その理由は粒子正孔対称性である。

\section{粒子正孔変換}
二部格子(すべての結合が 2 つの副格子 A、B をつなぐ格子)で、副格子の符号 $\varepsilon_i=\pm1$($i\in A$ で $+1$、$i\in B$ で $-1$)を使って
\begin{equation}
  \Pi:\quad c_{i\sigma}\to\varepsilon_i\cd_{i\sigma},\qquad \cd_{i\sigma}\to\varepsilon_ic_{i\sigma}
\end{equation}
と変換する。反交換関係を保つユニタリーな変換である。各項は次のように変わる。
\begin{itemize}
  \item 電子数:$n_{i\sigma}\to1-n_{i\sigma}$。
  \item ホッピング:結合の両端は副格子が違うので $\varepsilon_i\varepsilon_j=-1$。反交換の符号と合わせて $\cd_{i\sigma}c_{j\sigma}\to\cd_{j\sigma}c_{i\sigma}$ となり、運動エネルギーの項は全体として不変。二部格子でないと($t'$ や奇数の環)、ホッピングの符号が変わってしまう。
  \item 相互作用:$Un_{i\up}n_{i\dn}\to U(1-n_{i\up})(1-n_{i\dn})=Un_{i\up}n_{i\dn}-Un_i+U$。
\end{itemize}
まとめると $\Pi^\dagger H\Pi=H-U(\hat N-N_{\rm site})$ で、半充填のセクター $\hat N=N_{\rm site}$ の中では $H$ は不変である。$\Pi$ は $(N_\up,N_\dn)$ を $(N_{\rm site}-N_\up,N_{\rm site}-N_\dn)$ に写すので、半充填 $N_\up=N_\dn=N_{\rm site}/2$ のセクターは自分自身に写る。

\section{基底状態の一意性}
Lieb の定理\cite{lieb1989}によれば、有限で連結な二部格子上のハバード模型($U>0$、電子数が偶数)の基底状態は、スピンの多重項の縮退を除いて一意で、全スピンは $S=\lvert N_A-N_B\rvert/2$ である。$N_A=N_B$ なら完全に一意な一重項になる。本研究の二部格子の系はすべてこの条件を満たす。厳密対角化でも、1 次元開放端($U=4$、$L=4$〜$10$)で基底状態が一意で最初の励起が三重項であること、$U=0$ では縮退が残ることを確かめた。

一意な基底状態は $\Pi$ で自分自身に写り、位相しか付かない:$\Pi\ket{\psi_0}=e^{i\varphi}\ket{\psi_0}$。

\section{証明}
\begin{equation}
  \ev{n_{i\sigma}}=\bra{\psi_0}\Pi^\dagger n_{i\sigma}\Pi\ket{\psi_0}=\bra{\psi_0}(1-n_{i\sigma})\ket{\psi_0}=1-\ev{n_{i\sigma}}
\end{equation}
より $\ev{n_{i\sigma}}=1/2$、$\ev{n_i}=1$ が、すべてのサイトで成り立つ。サイトの位置を使っていないので、境界条件にも $U$ にもよらない。一般に、$\Pi$ で符号を変える演算子(電荷の偏り $n_i-1$、$Z_{i\up}$、$S^z_i$)の期待値は 0 で、符号を変えない演算子(オンサイト $ZZ$、$\Sv_i\cdot\Sv_j$、ホッピング)には制約がない。

\section{開放端の影響の行き先}
半充填での Friedel 振動の波数は $2k_F=\pi$ で、密度の交替 $(-1)^i$ に当たる。これはちょうど上の対称性で禁止される成分である。そのため開放端の影響は密度には現れず、対称性で禁止されない結合の強弱の交替(隣接スピン相関やホッピングの交替)として現れる。

\section{磁化を持つ UHF でも $\ev{n_i}=1$ になる理由}
$\Pi$ はスピンに対して $y$ 軸まわりの $\pi$ 回転として働き、$S^z_i$ の符号を変える。したがって $z$ 方向に磁化した UHF は $\Pi$ の対称性を破る。しかしハバード模型はスピン回転の対称性も持つので、組み合わせ $\tilde\Pi=\Pi R_y(\pi)^{-1}$ も対称性になる。$\tilde\Pi$ ではスピンは変わらず、電荷だけが反転する。磁化した UHF は $\tilde\Pi$ で自分自身に写るので、$\ev{n_i}=1$ を保つ。実際、二部格子の UHF では大きな磁化($\lvert\ev{S^z_i}\rvert$ が 0.44〜0.49)を持ちながら、$\lvert n_i-1\rvert$ は $10^{-15}$ の程度だった。

二部格子でない $4\times3$ のトーラスでは、この対称性がないので UHF の電荷が 6.5\% 偏った。

% =====================================================================
\chapter{和の観測量の分散の導出}
\label{app:sumvar}
% =====================================================================

命題\ref{prop:sumvar}を導く。観測量を $O=\sum_kc_kP_k$、台を $A_k$、各項の文字を $a^{(k)}_q$($q\in A_k$)とする。

\section{1 つの設定での推定値}
設定 $u$ で $n_m$ 回測り、$t$ 番目のショットの結果を $s^{(t)}_q$ とする。項 $k$ の推定値は
\begin{equation}
  X_k(u)=3^{\lvert A_k\rvert}\,\mathbf 1_k(u)\,\bar m_k,\qquad
  \mathbf 1_k(u)=\mathbf 1[\forall q\in A_k:\ b_q=a^{(k)}_q],\qquad
  \bar m_k=\frac1{n_m}\sum_{t=1}^{n_m}\prod_{q\in A_k}s^{(t)}_q
\end{equation}
で、観測量の推定値は $X(u)=\sum_kc_kX_k(u)$ である。CRM では $X^{\rm CRM}(u)=\sum_kc_k3^{\lvert A_k\rvert}\mathbf 1_k(u)(\bar m_k-y_k)+\sum_kc_ky_k$ となる。

\section{2 次のモーメント}
$\mathbb E[X(u)^2]=\sum_{k,l}c_kc_l\,3^{\lvert A_k\rvert+\lvert A_l\rvert}\,\mathbb E\bigl[\mathbf 1_k\mathbf 1_l\,\bar m_k\bar m_l\bigr]$ を計算する。
\begin{itemize}
  \item \textbf{基底の一致.} $\mathbf 1_k\mathbf 1_l=1$ となるのは、$A_k\cup A_l$ のすべての量子ビットで基底が指定どおりのときである。$A_k\cap A_l$ の上で $P_k$ と $P_l$ の文字が違えば(両立しなければ)ありえない。両立するとき、その確率は $3^{-\lvert A_k\cup A_l\rvert}$ である。$\lvert A_k\rvert+\lvert A_l\rvert-\lvert A_k\cup A_l\rvert=\lvert A_k\cap A_l\rvert$ なので、前の因子と合わせて $3^{\lvert A_k\cap A_l\rvert}$ が残る。
  \item \textbf{ショットの平均.} 基底が両方に一致した設定では、1 ショットの $\prod_{q\in A_k}s_q$ と $\prod_{q\in A_l}s_q$ は $P_k$ と $P_l$ の同時固有値で、その積は $P_kP_l$ の固有値である(両立する項は可換で、積は共通の量子ビットで恒等になる)。同じショットどうしの積の平均は $\ev{P_kP_l}$、別々のショットどうしは独立で $x_kx_l$ なので
  \begin{equation}
    \mathbb E[\bar m_k\bar m_l\mid\text{一致}]=\frac1{n_m}\ev{P_kP_l}+\Bigl(1-\frac1{n_m}\Bigr)x_kx_l.
  \end{equation}
\end{itemize}
したがって
\begin{equation}
  \mathbb E[X(u)^2]=\sum_{k,l\ \text{両立}}c_kc_l\,3^{\lvert A_k\cap A_l\rvert}\Bigl[\frac{\ev{P_kP_l}}{n_m}+\Bigl(1-\frac1{n_m}\Bigr)x_kx_l\Bigr]
\end{equation}
で、平均 $\mathbb E[X(u)]=\sum_kc_kx_k$ の 2 乗を引くと式\eqref{eq:sumvar-std}になる。

\section{CRM}
CRM では $\bar m_k$ を $\bar m_k-y_k$ に置き換える($y_k$ は定数)。一致した設定で
\begin{equation}
  \mathbb E[(\bar m_k-y_k)(\bar m_l-y_l)\mid\text{一致}]=\frac{\ev{P_kP_l}-x_kx_l}{n_m}+\Delta_k\Delta_l
\end{equation}
となる(ショットの共分散 $(\ev{P_kP_l}-x_kx_l)/n_m$ と、平均のずれ $\Delta_k\Delta_l$ の和)。定数項 $\sum_kc_ky_k$ は分散に寄与しないので、平均 $\sum_kc_k\Delta_k$ の 2 乗を引いて式\eqref{eq:sumvar-crm}を得る。

単一の項 $P$ では $\ev{P^2}=1$、$3^{\lvert A\cap A\rvert}=3^{\lvert A\rvert}$ なので、式\eqref{eq:var-std}・\eqref{eq:var-crm}に戻る。

% =====================================================================
\chapter{計算の再現}
\label{app:repro}
% =====================================================================

本論文の数値を出したスクリプトとデータは、リポジトリの \texttt{examples/tatsuyama/} にある。重い計算(DMRG、大規模なサンプリング)は研究室の計算機クラスター(PBS)で実行した。

\begin{center}\small
\begin{tabular}{@{}p{0.34\linewidth}p{0.58\linewidth}@{}}
\toprule
内容 & スクリプト / データ \\
\midrule
和の観測量の厳密な分散 & \texttt{crm\_pauli\_sum\_variance.py}、\texttt{crm\_exact\_gain.jl} \\
厳密な参照の 80 系の総当たり & \texttt{crm\_2d\_ed\_fid\_table.jl} / \texttt{crm\_edfid\_all.tsv} \\
系の大きさと利得(表\ref{tab:size}) & \texttt{crm\_mps\_scaling.jl} / \texttt{crm\_mps\_scaling\_results.tsv} \\
PHF・MPS の比較 & \texttt{crm\_phf.py}、\texttt{crm\_phf\_compare.py}、\texttt{crm\_mps\_lowchi.jl} \\
変分 MPS の回転平均 & \texttt{crm\_mps\_lowchi\_sym.jl}、\texttt{crm\_varsym\_compare.py} \\
2 次元 $W=4$ & \texttt{crm\_2d\_symuhf.jl} / \texttt{crm\_2d\_symuhf\_results.tsv} \\
2 次元 $W=6$、ドープ系 & \texttt{crm\_2d\_w6.jl}、\texttt{crm\_2d\_doped.jl}、\texttt{crm\_2d\_w6\_doped.jl} \\
相図とショット配分 & \texttt{crm\_param\_sweep.jl} / \texttt{crm\_sweep\_U.tsv}、\texttt{crm\_sweep\_nm.tsv} \\
読み出し誤差 & \texttt{crm\_readout\_noise.py} / \texttt{crm\_readout\_noise.out} \\
回転平均した prior の距離依存 & \texttt{crm\_varsym\_pairs.jl}、\texttt{crm\_varsym\_distance.py} \\
図\ref{fig:locality} & \texttt{crm\_fig\_locality\_exact.py} \\
粒子正孔対称性の確認 & \texttt{crm\_particle\_hole\_check.py} \\
\bottomrule
\end{tabular}
\end{center}

\section*{ソフトウェアの版}
\begin{center}\small
\begin{tabular}{@{}lll@{}}
\toprule
用途 & ソフトウェア & 版 \\
\midrule
DMRG、MPS、測定の模擬(クラスター) & Julia & 1.12.1 \\
 & ITensors.jl / ITensorMPS.jl & 0.9.30 / 0.4.1 \\
 & NDTensors.jl / KrylovKit.jl & 0.4.28 / 0.10.4 \\
平均場、表の集計(主にクラスター) & Python / NumPy & 3.13.13 / 2.5.1 \\
図(手元) & Python / NumPy / Matplotlib & 3.13.5 / 2.1.3 / 3.10.0 \\
本文 & Lua\LaTeX(TeX Live 2025) & LuaHBTeX 1.21.0 \\
\bottomrule
\end{tabular}
\end{center}
Julia のパッケージの版は、クラスター上の環境 \texttt{Hubbard\_MPS\_Env\_v2} の \texttt{Manifest.toml} による。


\begin{thebibliography}{99}
\bibitem{vermersch2024} B.~Vermersch, A.~Rath, B.~Sundar, C.~Branciard, J.~Preskill, and A.~Elben, ``Enhanced estimation of quantum properties with common randomized measurements,'' PRX Quantum \textbf{5}, 010352 (2024); arXiv:2304.12292.
\bibitem{huang2020} H.-Y.~Huang, R.~Kueng, and J.~Preskill, ``Predicting many properties of a quantum system from very few measurements,'' Nat.\ Phys.\ \textbf{16}, 1050 (2020).
\bibitem{huang2021} H.-Y.~Huang, R.~Kueng, and J.~Preskill, ``Efficient estimation of Pauli observables by derandomization,'' Phys.\ Rev.\ Lett.\ \textbf{127}, 030503 (2021).
\bibitem{zhao2021} A.~Zhao, N.~C.~Rubin, and A.~Miyake, ``Fermionic partial tomography via classical shadows,'' Phys.\ Rev.\ Lett.\ \textbf{127}, 110504 (2021).
\bibitem{wan2023} K.~Wan, W.~J.~Huggins, J.~Lee, and R.~Babbush, ``Matchgate shadows for fermionic quantum simulation,'' Commun.\ Math.\ Phys.\ \textbf{404}, 629 (2023).
\bibitem{lieb1989} E.~H.~Lieb, ``Two theorems on the Hubbard model,'' Phys.\ Rev.\ Lett.\ \textbf{62}, 1201 (1989).
\bibitem{itensor} M.~Fishman, S.~R.~White, and E.~M.~Stoudenmire, ``The ITensor software library for tensor network calculations,'' SciPost Phys.\ Codebases 4 (2022).
\end{thebibliography}
\todo{参考文献の書誌情報(巻・ページ)を原典で確認する。Hubbard 模型、DMRG、Hartree--Fock、スピン射影の標準的な文献を加える}


\end{document}
